% Abstract (~230 words). REWRITE 09-14 (사용자 결정): 구조(pathway) 위주 — 세 성질 drop-in /
% retrofittable / monotone-beyond-demos. π0.5 미언급 (사용자). 수치 출처 RESULTS_LINEUP + EVIDENCE §3.5, PI0 doc.
Vision-language-action (VLA) policies increasingly receive wrist wrench and contact signals in
their state, yet a behavior-cloned policy given the full wrench can still fail to bind its
behavior to it. On a real-robot suction case-picking task where the cases carry battery
cells---a $15$\,N contact-force limit, and a stop-on-contact transition observable only through
force---a naively fine-tuned VLA presses through contact in every trial and never stops on its
own. We present a conditioning pathway that changes how force is delivered, not what is
delivered: a condition vector \chat{} (contact, force magnitude, seal) computed from signals the
policy already receives, injected by FiLM at the state token, with the raw wrench masked. The
pathway has three properties. It is \emph{drop-in}: no backbone change, no new information, no
cost in imitation error, and it attaches to a second backbone ($\pi_0$). It is
\emph{retrofittable}: added to an already-trained naive policy, a few thousand steps of
continued training suffice. And it \emph{extrapolates in the right direction}: because \chat{} is a
calibrated monotone scalar, the policy brakes harder the harder it is pressed---stopping and
retreating at forces beyond the demonstrations, where the naive policy's response fades---and
picks succeed at a stack height absent from training. A matched ablation shows both
ingredients are necessary: unmasked, the pathway trains causally dead (\emph{bypass});
ungrounded, it is the worst policy in the study. We also ran it on our robot: every
conditioned pick stops on its own (median contact force $2.5$\,N) where the naive policy is
externally aborted in every attempt, and on-robot counterfactual probes reproduce the offline
crossover.
